\clearpage
\chapter*{Erratum à l'exposé n\textsuperscript{o} 221}
\addcontentsline{toc}{chapter}{Erratum à l'exposé n\textsuperscript{o} 221}
\setcounter{footnote}{3}
\src{302}

\noindent\textbf{6. TDTE IV : Les schémas de Hilbert.}

\begin{center}
[Séminaire Bourbaki, t. 13, 1960/61, n\textsuperscript{o} 221, 28 p.]
\end{center}

\noindent\textbf{Page 6, théorème 2.1, 3\textsuperscript{e} ligne.} --- Au lieu
de « il faut », lire « il faut et il suffit ».

\noindent\textbf{Page 15, ligne 13.} --- Au lieu de « $P\neq Q$ », lire
« $P<Q$ ».

\noindent\textbf{Ligne 15.} --- Au lieu de « Si au contraire $P=Q$ », lire
« Dans le cas contraire, on aura $P(n)\geq Q(n)$ pour $n$ grand ».

\noindent\textbf{Ligne 18, formule (*).} --- Au lieu de « $<$ » et « $>$ »,
lire « $\leq$ » et « $\geq$ ».

\noindent\textbf{Ligne 5 à partir du bas.} --- Au lieu de « $r$ », lire
« $Z$ ».

\noindent\textbf{Page 18, remarque 3.9.} --- Signalons que l'étude des
composantes connexes des schémas de Hilbert sur un corps algébriquement clos a
été faite par HARTSHORNE. Cet auteur prouve que les
$\operatorname{Hilb}^{P}_{r}$ sont connexes, et détermine les couples $(r,P)$
pour lesquels $\operatorname{Hilb}^{P}_{r}\neq\varnothing$\footnote{HARTSHORNE
(R.). --- Connectedness of the Hilbert scheme (Thesis, Harvard, 1962).}.

\noindent\textbf{Page 23, ligne 5 à partir du bas.} --- Au lieu de « une
composante », lire « d'une composante ».

\noindent\textbf{Pages 23 et 24, remarque 5.5.} --- À propos de l'exemple de
ZAPPA, signalons que MUMFORD vient même de construire une composante
irréductible du schéma de Hilbert pour $\mathbf P^3_{\mathbf C}$ (dont les
points généraux représentent des courbes non singulières de degré 14, genre
24), qui est non réduite en son point générique. Faisant éclater une des
courbes obtenues, il obtient également un schéma projectif régulier de
dimension 3 sur $\mathbf C$, dont le schéma formel des modules est non réduit
en son point générique, ou, ce qui revient au même, tel que sa variété des
modules locale, au sens de la géométrie analytique, sur $\mathbf C$ (cf.
Séminaire Cartan, t. 13, 1960/61), est non réduite en tous ses points.
